\section{Conclusion}\label{sec:conclusion}

A quantum slice helps a developer only if it never hides code that can change the qubits asked about. We defined that
requirement as preserving the qubits' joint state for every input while keeping every semantic dependence, and showed
that the write-all rule meets it. Its slice is the backward light cone extended along classical bits, needs only
classical data, classical control and quantum data dependence, and is computed in one pass. A pre-registered study
measured what the other rules lose. Reading a control as unchanged gives wrong slices that often reproduce the
program's outcomes, entanglement edges add only code that cannot matter, and pruning trades dependences for size and
holds for one observable or input. The QDG slice comes within a few points of the smallest complete slice and costs
milliseconds. Two questions remain, whether these slices shorten the manual reading that developers report does not
scale, which needs a study with developers, and how to compute them statically on the source developers write, at
the scale of real programs and with loops that run on the device.
